\section{AI Agent 的可观测性与安全控制}

\subsection{监控指标与报警策略}

一个成熟的 Agent 平台首先要定义可观测指标，然后在「信号」出现时自动告警：

\begin{itemize}
    \item \textbf{语义一致性}：Historical Reference vs. current response divergence
    \item \textbf{模型幻觉率}：随机抽查输出，看是否引入虚假事实
    \item \textbf{工具调用错误}：命令执行失败 / API 调用超时
    \item \textbf{上下文暴涨}：上下文窗口突然膨胀说明 memory leak 或 prompt injection
    \item \textbf{用户纠错次数}：用户频繁打断说明 agent 误解 intent
\end{itemize}

\begin{knowledgebox}{越早捕捉 signals 越容易纠偏}
很多团队等到用户投诉才知道 agent 出问题了。早期的异常信号往往隐藏在 latency spike、工具调用超时或“内容重复”中。设置透明的「metric + 增强 sampling + human-in-loop review」可在 agent 出错前升级响应。
\end{knowledgebox}

Burak 强调 instrumentation 的数据 pipeline 要保持低开销：对语义一致性或 hallucination 只采样 1-2\% 请求，通过 squad 评审保持敏感性。过度采集会产生成本，也可能导致告警疲劳。

\subsection{安全治理与 guardrail 设计}

Agent 的行为不能只靠模型 prompt 设定，还要靠系统级 guardrail：

\begin{itemize}
    \item \textbf{Policy Layer}：在调用前拦截敏感操作（删除文件、调度任务）
    \item \textbf{Tool Firewall}：将模型对工具的访问限制在白名单
    \item \textbf{Human-in-the-loop level}：根据任务风险动态决定是否需要人工审批
    \item \textbf{Audit Trail}：记录每一步决策、工具调用和用户确认
\end{itemize}

\begin{warningbox}{没有 audit trail 就无法追责}
一旦 agent 造成错误或安全事件，缺少操作日志意味着无法恢复决策链。这会导致治理系统失效，甚至影响合规审计。Anthropic 和 Vertex AI 都强调：每一步 tool call 都要可追溯、可还原。
\end{warningbox}

\subsection{本章小结}

通过量化指标、报警链路和责任链路，企业才能在 agent 出现异常时迅速反应，并在能力上升的同时保持安全边界。

